{
\captionsetup{type=figure}
\captionof{figure}{Trajectory for instance \texttt{django\_\_django-14672}.}
\begin{issuebox}[django\_\_django-14672]We're currently solving the following issue within our repository. Here's the issue text:

ISSUE:

Missing call `make\_hashable' on `through\_fields' in `ManyToManyRel'

Description

	

In 3.2 identity property has been added to all ForeignObjectRel to make it possible to compare them. A hash is derived from said identity and it's possible because identity is a tuple. To make limit\_choices\_to hashable (one of this tuple elements), there's a call to make\_hashable.

It happens that through\_fields can be a list. In such case, this make\_hashable call is missing in ManyToManyRel.

For some reason it only fails on checking proxy model. I think proxy models have 29 checks and normal ones 24, hence the issue, but that's just a guess.

Minimal repro:
\begin{CodeVerbatim}[breaklines=true, breakanywhere=true]
class Parent(models.Model):
	name = models.CharField(max\_length=256)
class ProxyParent(Parent):
	class Meta:
		proxy = True
class Child(models.Model):
	parent = models.ForeignKey(Parent, on\_delete=models.CASCADE)
	many\_to\_many\_field = models.ManyToManyField(
		to=Parent,
		through="ManyToManyModel",
		through\_fields=['child', 'parent'],
		related\_name="something"
	)

class ManyToManyModel(models.Model):
	parent = models.ForeignKey(Parent, on\_delete=models.CASCADE, related\_name='+')
	child = models.ForeignKey(Child, on\_delete=models.CASCADE, related\_name='+')
	second\_child = models.ForeignKey(Child, on\_delete=models.CASCADE, null=True, default=None)
\end{CodeVerbatim}

Which will result in 
\begin{CodeVerbatim}[breaklines=true, breakanywhere=true]
 File "manage.py", line 23, in <module>
	main()
 File "manage.py", line 19, in main
	execute\_from\_command\_line(sys.argv)

 File "/home/tom/PycharmProjects/broken\_m2m\_project/venv/lib/python3.8/site-packages/django/core/management/\_\_init\_\_.py", line 419, in execute\_from\_command\_line
	utility.execute()

 File "/home/tom/PycharmProjects/broken\_m2m\_project/venv/lib/python3.8/site-packages/django/core/management/\_\_init\_\_.py", line 413, in execute
	self.fetch\_command(subcommand).run\_from\_argv(self.argv)

 File "/home/tom/PycharmProjects/broken\_m2m\_project/venv/lib/python3.8/site-packages/django/core/management/base.py", line 354, in run\_from\_argv
	self.execute(*args, **cmd\_options)

 File "/home/tom/PycharmProjects/broken\_m2m\_project/venv/lib/python3.8/site-packages/django/core/management/base.py", line 393, in execute
	self.check()

 File "/home/tom/PycharmProjects/broken\_m2m\_project/venv/lib/python3.8/site-packages/django/core/management/base.py", line 419, in check
	all\_issues = checks.run\_checks(

 File "/home/tom/PycharmProjects/broken\_m2m\_project/venv/lib/python3.8/site-packages/django/core/checks/registry.py", line 76, in run\_checks
	new\_errors = check(app\_configs=app\_configs, databases=databases)

 File "/home/tom/PycharmProjects/broken\_m2m\_project/venv/lib/python3.8/site-packages/django/core/checks/model\_checks.py", line 34, in check\_all\_models
	errors.extend(model.check(**kwargs))

 File "/home/tom/PycharmProjects/broken\_m2m\_project/venv/lib/python3.8/site-packages/django/db/models/base.py", line 1277, in check
	*cls.\_check\_field\_name\_clashes(),

 File "/home/tom/PycharmProjects/djangbroken\_m2m\_projectProject/venv/lib/python3.8/site-packages/django/db/models/base.py", line 1465, in \_check\_field\_name\_clashes
	if f not in used\_fields:

 File "/home/tom/PycharmProjects/broken\_m2m\_project/venv/lib/python3.8/site-packages/django/db/models/fields/reverse\_related.py", line 140, in \_\_hash\_\_
	return hash(self.identity)

TypeError: unhashable type: 'list'
\end{CodeVerbatim}

Solution: Add missing make\_hashable call on self.through\_fields in ManyToManyRel.


INSTRUCTIONS: (31 lines)

\begin{CodeVerbatim}[breaklines=true, breakanywhere=true]
(Open file: n/a)
(Current directory: /django__django)
bash-\$
\end{CodeVerbatim}
\end{issuebox}
\begin{agentbox}[django\_\_django-14672]
To address the issue, we should first replicate the bug by creating a script that mimics the minimal reproduction example provided in the issue description. This script will attempt to perform the operations that lead to the `TypeError: unhashable type: 'list'` error. Let's start by creating a new Python script named `reproduce.py`.



\hrulefill
\begin{CodeVerbatim}[breaklines=true, breakanywhere=true]
create reproduce.py
\end{CodeVerbatim}
\end{agentbox}
\begin{observationbox}[django\_\_django-14672]\begin{Verbatim}[breaklines=true, breakanywhere=true]
[File: /django__django/reproduce.py (1 lines total)]
1:

(Open file: n/a)
(Current directory: /django__django)
bash-$
\end{Verbatim}
\end{observationbox}
\begin{agentbox}[django\_\_django-14672]
Now let's paste in the minimal reproduction code from the issue into `reproduce.py`.



\hrulefill
\begin{CodeVerbatim}[breaklines=true, breakanywhere=true]
edit 1:1
from django.db import models

class Parent(models.Model):
    name = models.CharField(max_length=256)

class ProxyParent(Parent):
    class Meta:
        proxy = True

class Child(models.Model):
    parent = models.ForeignKey(Parent, on_delete=models.CASCADE)
    many_to_many_field = models.ManyToManyField(
        to=Parent,
        through="ManyToManyModel",
        through_fields=['child', 'parent'],
        related_name="something"
    )

class ManyToManyModel(models.Model):
    parent = models.ForeignKey(Parent, on_delete=models.CASCADE, related_name='+')
    child = models.ForeignKey(Child, on_delete=models.CASCADE, related_name='+')
    second_child = models.ForeignKey(Child, on_delete=models.CASCADE, null=True, default=None)

print("Script completed successfully, no errors.")
end_of_edit
\end{CodeVerbatim}
\end{agentbox}
\begin{observationbox}[django\_\_django-14672]\begin{Verbatim}[breaklines=true, breakanywhere=true]
[File: /django__django/reproduce.py (24 lines total)]
1:from django.db import models
2:
3:class Parent(models.Model):
4:    name = models.CharField(max_length=256)
5:
6:class ProxyParent(Parent):
7:    class Meta:
8:        proxy = True
9:
10:class Child(models.Model):
11:    parent = models.ForeignKey(Parent, on_delete=models.CASCADE)
12:    many_to_many_field = models.ManyToManyField(
13:        to=Parent,
14:        through="ManyToManyModel",
15:        through_fields=['child', 'parent'],
16:        related_name="something"
17:    )
18:
19:class ManyToManyModel(models.Model):
20:    parent = models.ForeignKey(Parent, on_delete=models.CASCADE, related_name='+')
21:    child = models.ForeignKey(Child, on_delete=models.CASCADE, related_name='+')
22:    second_child = models.ForeignKey(Child, on_delete=models.CASCADE, null=True, default=None)
23:
24:print("Script completed successfully, no errors.")
File updated. Please review the changes and make sure they are correct (correct indentation, no duplicate lines, etc). Edit the file again if necessary.

(Open file: /django__django/reproduce.py)
(Current directory: /django__django)
bash-$
\end{Verbatim}
\end{observationbox}
\begin{agentbox}[django\_\_django-14672]
The reproduction script has been created successfully with the correct indentation and content. Now, let's run the script to see if we can reproduce the bug.



\hrulefill
\begin{CodeVerbatim}[breaklines=true, breakanywhere=true]
python reproduce.py
\end{CodeVerbatim}
\end{agentbox}
\begin{observationbox}[django\_\_django-14672]\begin{Verbatim}[breaklines=true, breakanywhere=true]
Traceback (most recent call last):
  File "reproduce.py", line 3, in <module>
    class Parent(models.Model):
  File "/django__django/django/db/models/base.py", line 108, in __new__
    app_config = apps.get_containing_app_config(module)
  File "/django__django/django/apps/registry.py", line 253, in get_containing_app_config
    self.check_apps_ready()
  File "/django__django/django/apps/registry.py", line 135, in check_apps_ready
    settings.INSTALLED_APPS
  File "/django__django/django/conf/__init__.py", line 69, in __getattr__
    self._setup(name)
  File "/django__django/django/conf/__init__.py", line 50, in _setup
    raise ImproperlyConfigured(
django.core.exceptions.ImproperlyConfigured: Requested setting INSTALLED_APPS, but settings are not configured. You must either define the environment variable DJANGO_SETTINGS_MODULE or call settings.configure() before accessing settings.

(Open file: /django__django/reproduce.py)
(Current directory: /django__django)
bash-$
\end{Verbatim}
\end{observationbox}
\begin{agentbox}[django\_\_django-14672]
The error indicates that Django settings are not configured, which is expected because we're running a standalone script that uses Django models outside of a Django project environment. To fix this, we need to set up the Django environment within the script before defining the models.



Let's edit `reproduce.py` to configure the Django settings at the beginning of the script.



\hrulefill
\begin{CodeVerbatim}[breaklines=true, breakanywhere=true]
edit 1:1
import django
from django.conf import settings

settings.configure(
    INSTALLED_APPS=[
        'django.contrib.auth',
        'django.contrib.contenttypes',
        '__main__',  # This will include the models we define in this script
    ],
    DATABASES={
        'default': {
            'ENGINE': 'django.db.backends.sqlite3',
            'NAME': 'mydatabase',
        }
    }
)
django.setup()

from django.db import models

class Parent(models.Model):
    name = models.CharField(max_length=256)

class ProxyParent(Parent):
    class Meta:
        proxy = True

class Child(models.Model):
    parent = models.ForeignKey(Parent, on_delete=models.CASCADE)
    many_to_many_field = models.ManyToManyField(
        to=Parent,
        through="ManyToManyModel",
        through_fields=['child', 'parent'],
        related_name="something"
    )

class ManyToManyModel(models.Model):
    parent = models.ForeignKey(Parent, on_delete=models.CASCADE, related_name='+')
    child = models.ForeignKey(Child, on_delete=models.CASCADE, related_name='+')
    second_child = models.ForeignKey(Child, on_delete=models.CASCADE, null=True, default=None)

print("Script completed successfully, no errors.")
end_of_edit
\end{CodeVerbatim}
\end{agentbox}
\begin{observationbox}[django\_\_django-14672]\begin{Verbatim}[breaklines=true, breakanywhere=true]
[File: /django__django/reproduce.py (65 lines total)]
1:import django
2:from django.conf import settings
3:
4:settings.configure(
5:    INSTALLED_APPS=[
6:        'django.contrib.auth',
7:        'django.contrib.contenttypes',
8:        '__main__',  # This will include the models we define in this script
9:    ],
10:    DATABASES={
11:        'default': {
12:            'ENGINE': 'django.db.backends.sqlite3',
13:            'NAME': 'mydatabase',
14:        }
15:    }
16:)
17:django.setup()
18:
19:from django.db import models
20:
21:class Parent(models.Model):
22:    name = models.CharField(max_length=256)
23:
24:class ProxyParent(Parent):
25:    class Meta:
26:        proxy = True
27:
28:class Child(models.Model):
29:    parent = models.ForeignKey(Parent, on_delete=models.CASCADE)
30:    many_to_many_field = models.ManyToManyField(
31:        to=Parent,
32:        through="ManyToManyModel",
33:        through_fields=['child', 'parent'],
34:        related_name="something"
35:    )
36:
37:class ManyToManyModel(models.Model):
38:    parent = models.ForeignKey(Parent, on_delete=models.CASCADE, related_name='+')
39:    child = models.ForeignKey(Child, on_delete=models.CASCADE, related_name='+')
40:    second_child = models.ForeignKey(Child, on_delete=models.CASCADE, null=True, default=None)
41:
42:print("Script completed successfully, no errors.")
43:
44:class Parent(models.Model):
45:    name = models.CharField(max_length=256)
46:
47:class ProxyParent(Parent):
48:    class Meta:
49:        proxy = True
50:
51:class Child(models.Model):
52:    parent = models.ForeignKey(Parent, on_delete=models.CASCADE)
53:    many_to_many_field = models.ManyToManyField(
54:        to=Parent,
55:        through="ManyToManyModel",
56:        through_fields=['child', 'parent'],
57:        related_name="something"
58:    )
59:
60:class ManyToManyModel(models.Model):
61:    parent = models.ForeignKey(Parent, on_delete=models.CASCADE, related_name='+')
62:    child = models.ForeignKey(Child, on_delete=models.CASCADE, related_name='+')
63:    second_child = models.ForeignKey(Child, on_delete=models.CASCADE, null=True, default=None)
64:
65:print("Script completed successfully, no errors.")
File updated. Please review the changes and make sure they are correct (correct indentation, no duplicate lines, etc). Edit the file again if necessary.

(Open file: /django__django/reproduce.py)
(Current directory: /django__django)
bash-$
\end{Verbatim}
\end{observationbox}
\begin{agentbox}[django\_\_django-14672]
It seems that the models have been duplicated in the `reproduce.py` file. This is likely an error from the edit command. We need to remove the duplicate model definitions to ensure the script is correct.



Let's edit the file again to remove the duplicate lines.



\hrulefill
\begin{CodeVerbatim}[breaklines=true, breakanywhere=true]
edit 44:65
end_of_edit
\end{CodeVerbatim}
\end{agentbox}
\begin{observationbox}[django\_\_django-14672]\begin{Verbatim}[breaklines=true, breakanywhere=true]
[File: /django__django/reproduce.py (43 lines total)]
1:import django
2:from django.conf import settings
3:
4:settings.configure(
5:    INSTALLED_APPS=[
6:        'django.contrib.auth',
7:        'django.contrib.contenttypes',
8:        '__main__',  # This will include the models we define in this script
9:    ],
10:    DATABASES={
11:        'default': {
12:            'ENGINE': 'django.db.backends.sqlite3',
13:            'NAME': 'mydatabase',
14:        }
15:    }
16:)
17:django.setup()
18:
19:from django.db import models
20:
21:class Parent(models.Model):
22:    name = models.CharField(max_length=256)
23:
24:class ProxyParent(Parent):
25:    class Meta:
26:        proxy = True
27:
28:class Child(models.Model):
29:    parent = models.ForeignKey(Parent, on_delete=models.CASCADE)
30:    many_to_many_field = models.ManyToManyField(
31:        to=Parent,
32:        through="ManyToManyModel",
33:        through_fields=['child', 'parent'],
34:        related_name="something"
35:    )
36:
37:class ManyToManyModel(models.Model):
38:    parent = models.ForeignKey(Parent, on_delete=models.CASCADE, related_name='+')
39:    child = models.ForeignKey(Child, on_delete=models.CASCADE, related_name='+')
40:    second_child = models.ForeignKey(Child, on_delete=models.CASCADE, null=True, default=None)
41:
42:print("Script completed successfully, no errors.")
43:
File updated. Please review the changes and make sure they are correct (correct indentation, no duplicate lines, etc). Edit the file again if necessary.

(Open file: /django__django/reproduce.py)
(Current directory: /django__django)
bash-$
\end{Verbatim}
\end{observationbox}
\begin{agentbox}[django\_\_django-14672]
The duplicate lines have been successfully removed from the `reproduce.py` file. Now, let's run the script again to see if we can reproduce the bug.



\hrulefill
\begin{CodeVerbatim}[breaklines=true, breakanywhere=true]
python reproduce.py
\end{CodeVerbatim}
\end{agentbox}
\begin{observationbox}[django\_\_django-14672]\begin{Verbatim}[breaklines=true, breakanywhere=true]
Script completed successfully, no errors.

(Open file: /django__django/reproduce.py)
(Current directory: /django__django)
bash-$
\end{Verbatim}
\end{observationbox}
\begin{agentbox}[django\_\_django-14672]
The script ran without errors, which means we haven't successfully reproduced the bug yet. The issue description mentions that the error occurs during a check, which is likely part of Django's system checks framework. We need to trigger these checks to reproduce the error.



Let's modify the `reproduce.py` script to include a call to Django's check framework after setting up the models.



\hrulefill
\begin{CodeVerbatim}[breaklines=true, breakanywhere=true]
edit 42:42
from django.core.management import call_command

# Trigger Django's check framework
call_command('check')
end_of_edit
\end{CodeVerbatim}
\end{agentbox}
\begin{observationbox}[django\_\_django-14672]\begin{Verbatim}[breaklines=true, breakanywhere=true]
[File: /django__django/reproduce.py (46 lines total)]
1:import django
2:from django.conf import settings
3:
4:settings.configure(
5:    INSTALLED_APPS=[
6:        'django.contrib.auth',
7:        'django.contrib.contenttypes',
8:        '__main__',  # This will include the models we define in this script
9:    ],
10:    DATABASES={
11:        'default': {
12:            'ENGINE': 'django.db.backends.sqlite3',
13:            'NAME': 'mydatabase',
14:        }
15:    }
16:)
17:django.setup()
18:
19:from django.db import models
20:
21:class Parent(models.Model):
22:    name = models.CharField(max_length=256)
23:
24:class ProxyParent(Parent):
25:    class Meta:
26:        proxy = True
27:
28:class Child(models.Model):
29:    parent = models.ForeignKey(Parent, on_delete=models.CASCADE)
30:    many_to_many_field = models.ManyToManyField(
31:        to=Parent,
32:        through="ManyToManyModel",
33:        through_fields=['child', 'parent'],
34:        related_name="something"
35:    )
36:
37:class ManyToManyModel(models.Model):
38:    parent = models.ForeignKey(Parent, on_delete=models.CASCADE, related_name='+')
39:    child = models.ForeignKey(Child, on_delete=models.CASCADE, related_name='+')
40:    second_child = models.ForeignKey(Child, on_delete=models.CASCADE, null=True, default=None)
41:
42:from django.core.management import call_command
43:
44:# Trigger Django's check framework
45:call_command('check')
46:
File updated. Please review the changes and make sure they are correct (correct indentation, no duplicate lines, etc). Edit the file again if necessary.

(Open file: /django__django/reproduce.py)
(Current directory: /django__django)
bash-$
\end{Verbatim}
\end{observationbox}
\begin{agentbox}[django\_\_django-14672]
The `reproduce.py` script has been updated to include a call to Django's check framework. Let's run the script again to see if we can reproduce the bug.



\hrulefill
\begin{CodeVerbatim}[breaklines=true, breakanywhere=true]
python reproduce.py
\end{CodeVerbatim}
\end{agentbox}
\begin{observationbox}[django\_\_django-14672]\begin{Verbatim}[breaklines=true, breakanywhere=true]
Traceback (most recent call last):
  File "reproduce.py", line 45, in <module>
    call_command('check')
  File "/django__django/django/core/management/__init__.py", line 181, in call_command
    return command.execute(*args, **defaults)
  File "/django__django/django/core/management/base.py", line 398, in execute
    output = self.handle(*args, **options)
  File "/django__django/django/core/management/commands/check.py", line 63, in handle
    self.check(
  File "/django__django/django/core/management/base.py", line 419, in check
    all_issues = checks.run_checks(
  File "/django__django/django/core/checks/registry.py", line 77, in run_checks
    new_errors = check(app_configs=app_configs, databases=databases)
  File "/django__django/django/core/checks/model_checks.py", line 34, in check_all_models
    errors.extend(model.check(**kwargs))
  File "/django__django/django/db/models/base.py", line 1281, in check
    *cls._check_field_name_clashes(),
  File "/django__django/django/db/models/base.py", line 1471, in _check_field_name_clashes
    if f not in used_fields:
  File "/django__django/django/db/models/fields/reverse_related.py", line 139, in __hash__
    return hash(self.identity)
TypeError: unhashable type: 'list'

(Open file: /django__django/reproduce.py)
(Current directory: /django__django)
bash-$
\end{Verbatim}
\end{observationbox}
\begin{agentbox}[django\_\_django-14672]
We have successfully reproduced the bug, which is confirmed by the `TypeError: unhashable type: 'list'` in the traceback. According to the issue description, the solution is to add a missing call to `make\_hashable` on `self.through\_fields` in `ManyToManyRel`.



To fix this, we need to locate the `ManyToManyRel` class within the Django source code and add the `make\_hashable` call where `self.through\_fields` is used in the `identity` property.



First, we need to find the file that contains the `ManyToManyRel` class. Since we're working within the Django repository, we can search for the file directly.



\hrulefill
\begin{CodeVerbatim}[breaklines=true, breakanywhere=true]
search_dir "class ManyToManyRel"
\end{CodeVerbatim}
\end{agentbox}
\begin{observationbox}[django\_\_django-14672]\begin{Verbatim}[breaklines=true, breakanywhere=true]
Found 2 matches for "class ManyToManyRel" in /django__django:
/django__django/django/db/models/fields/reverse_related.py (1 matches)
/django__django/tests/invalid_models_tests/test_relative_fields.py (1 matches)
End of matches for "class ManyToManyRel" in /django__django

(Open file: /django__django/reproduce.py)
(Current directory: /django__django)
bash-$
\end{Verbatim}
\end{observationbox}
\begin{agentbox}[django\_\_django-14672]
The `ManyToManyRel` class is found in the file `django/db/models/fields/reverse\_related.py`. Let's open this file to locate the exact place where we need to add the `make\_hashable` call.



\hrulefill
\begin{CodeVerbatim}[breaklines=true, breakanywhere=true]
open django/db/models/fields/reverse_related.py
\end{CodeVerbatim}
\end{agentbox}
\begin{observationbox}[django\_\_django-14672]\begin{Verbatim}[breaklines=true, breakanywhere=true]
[File: /django__django/django/db/models/fields/reverse_related.py (330 lines total)]
1:"""
2:"Rel objects" for related fields.
3:
4:"Rel objects" (for lack of a better name) carry information about the relation
5:modeled by a related field and provide some utility functions. They're stored
6:in the ``remote_field`` attribute of the field.
7:
8:They also act as reverse fields for the purposes of the Meta API because
9:they're the closest concept currently available.
10:"""
11:
12:from django.core import exceptions
13:from django.utils.functional import cached_property
14:from django.utils.hashable import make_hashable
15:
16:from . import BLANK_CHOICE_DASH
17:from .mixins import FieldCacheMixin
18:
19:
20:class ForeignObjectRel(FieldCacheMixin):
21:    """
22:    Used by ForeignObject to store information about the relation.
23:
24:    ``_meta.get_fields()`` returns this class to provide access to the field
25:    flags for the reverse relation.
26:    """
27:
28:    # Field flags
29:    auto_created = True
30:    concrete = False
31:    editable = False
32:    is_relation = True
33:
34:    # Reverse relations are always nullable (Django can't enforce that a
35:    # foreign key on the related model points to this model).
36:    null = True
37:    empty_strings_allowed = False
38:
39:    def __init__(self, field, to, related_name=None, related_query_name=None,
40:                 limit_choices_to=None, parent_link=False, on_delete=None):
41:        self.field = field
42:        self.model = to
43:        self.related_name = related_name
44:        self.related_query_name = related_query_name
45:        self.limit_choices_to = {} if limit_choices_to is None else limit_choices_to
46:        self.parent_link = parent_link
47:        self.on_delete = on_delete
48:
49:        self.symmetrical = False
50:        self.multiple = True
51:
52:    # Some of the following cached_properties can't be initialized in
53:    # __init__ as the field doesn't have its model yet. Calling these methods
54:    # before field.contribute_to_class() has been called will result in
55:    # AttributeError
56:    @cached_property
57:    def hidden(self):
58:        return self.is_hidden()
59:
60:    @cached_property
61:    def name(self):
62:        return self.field.related_query_name()
63:
64:    @property
65:    def remote_field(self):
66:        return self.field
67:
68:    @property
69:    def target_field(self):
70:        """
71:        When filtering against this relation, return the field on the remote
72:        model against which the filtering should happen.
73:        """
74:        target_fields = self.get_path_info()[-1].target_fields
75:        if len(target_fields) > 1:
76:            raise exceptions.FieldError("Can't use target_field for multicolumn relations.")
77:        return target_fields[0]
78:
79:    @cached_property
80:    def related_model(self):
81:        if not self.field.model:
82:            raise AttributeError(
83:                "This property can't be accessed before self.field.contribute_to_class has been called.")
84:        return self.field.model
85:
86:    @cached_property
87:    def many_to_many(self):
88:        return self.field.many_to_many
89:
90:    @cached_property
91:    def many_to_one(self):
92:        return self.field.one_to_many
93:
94:    @cached_property
95:    def one_to_many(self):
96:        return self.field.many_to_one
97:
98:    @cached_property
99:    def one_to_one(self):
100:        return self.field.one_to_one
(230 more lines below)

(Open file: /django__django/reproduce.py)
(Current directory: /django__django)
bash-$
\end{Verbatim}
\end{observationbox}
\begin{agentbox}[django\_\_django-14672]
The file `reverse\_related.py` is now open. We need to find the `identity` property within the `ManyToManyRel` class to add the `make\_hashable` call. Since the file is 330 lines long and we don't know the exact line number, we should search within this file for the `identity` property.



\hrulefill
\begin{CodeVerbatim}[breaklines=true, breakanywhere=true]
search_file "def identity(self)"
\end{CodeVerbatim}
\end{agentbox}
\begin{observationbox}[django\_\_django-14672]\begin{Verbatim}[breaklines=true, breakanywhere=true]
Found 3 matches for "def identity(self)" in /django__django/django/db/models/fields/reverse_related.py:
Line 120:    def identity(self):
Line 242:    def identity(self):
Line 310:    def identity(self):
End of matches for "def identity(self)" in /django__django/django/db/models/fields/reverse_related.py

(Open file: /django__django/django/db/models/fields/reverse_related.py)
(Current directory: /django__django)
bash-$
\end{Verbatim}
\end{observationbox}
\begin{agentbox}[django\_\_django-14672]
We have found three occurrences of the `identity` property in the `reverse\_related.py` file. Since we are interested in the `ManyToManyRel` class, we should navigate to the second occurrence, which is likely to be the correct one.



Let's use the `goto` command to move to the line just before the second occurrence of the `identity` property to inspect the context and confirm that it's within the `ManyToManyRel` class.



\hrulefill
\begin{CodeVerbatim}[breaklines=true, breakanywhere=true]
goto 241
\end{CodeVerbatim}
\end{agentbox}
\begin{observationbox}[django\_\_django-14672]\begin{Verbatim}[breaklines=true, breakanywhere=true]
[File: /django__django/django/db/models/fields/reverse_related.py (330 lines total)]
(225 more lines above)
226:            field, to,
227:            related_name=related_name,
228:            related_query_name=related_query_name,
229:            limit_choices_to=limit_choices_to,
230:            parent_link=parent_link,
231:            on_delete=on_delete,
232:        )
233:
234:        self.field_name = field_name
235:
236:    def __getstate__(self):
237:        state = self.__dict__.copy()
238:        state.pop('related_model', None)
239:        return state
240:
241:    @property
242:    def identity(self):
243:        return super().identity + (self.field_name,)
244:
245:    def get_related_field(self):
246:        """
247:        Return the Field in the 'to' object to which this relationship is tied.
248:        """
249:        field = self.model._meta.get_field(self.field_name)
250:        if not field.concrete:
251:            raise exceptions.FieldDoesNotExist("No related field named '%s'" % self.field_name)
252:        return field
253:
254:    def set_field_name(self):
255:        self.field_name = self.field_name or self.model._meta.pk.name
256:
257:
258:class OneToOneRel(ManyToOneRel):
259:    """
260:    Used by OneToOneField to store information about the relation.
261:
262:    ``_meta.get_fields()`` returns this class to provide access to the field
263:    flags for the reverse relation.
264:    """
265:
266:    def __init__(self, field, to, field_name, related_name=None, related_query_name=None,
267:                 limit_choices_to=None, parent_link=False, on_delete=None):
268:        super().__init__(
269:            field, to, field_name,
270:            related_name=related_name,
271:            related_query_name=related_query_name,
272:            limit_choices_to=limit_choices_to,
273:            parent_link=parent_link,
274:            on_delete=on_delete,
275:        )
276:
277:        self.multiple = False
278:
279:
280:class ManyToManyRel(ForeignObjectRel):
281:    """
282:    Used by ManyToManyField to store information about the relation.
283:
284:    ``_meta.get_fields()`` returns this class to provide access to the field
285:    flags for the reverse relation.
286:    """
287:
288:    def __init__(self, field, to, related_name=None, related_query_name=None,
289:                 limit_choices_to=None, symmetrical=True, through=None,
290:                 through_fields=None, db_constraint=True):
291:        super().__init__(
292:            field, to,
293:            related_name=related_name,
294:            related_query_name=related_query_name,
295:            limit_choices_to=limit_choices_to,
296:        )
297:
298:        if through and not db_constraint:
299:            raise ValueError("Can't supply a through model and db_constraint=False")
300:        self.through = through
301:
302:        if through_fields and not through:
303:            raise ValueError("Cannot specify through_fields without a through model")
304:        self.through_fields = through_fields
305:
306:        self.symmetrical = symmetrical
307:        self.db_constraint = db_constraint
308:
309:    @property
310:    def identity(self):
311:        return super().identity + (
312:            self.through,
313:            self.through_fields,
314:            self.db_constraint,
315:        )
316:
317:    def get_related_field(self):
318:        """
319:        Return the field in the 'to' object to which this relationship is tied.
320:        Provided for symmetry with ManyToOneRel.
321:        """
322:        opts = self.through._meta
323:        if self.through_fields:
324:            field = opts.get_field(self.through_fields[0])
325:        else:
(5 more lines below)

(Open file: /django__django/django/db/models/fields/reverse_related.py)
(Current directory: /django__django)
bash-$
\end{Verbatim}
\end{observationbox}
\begin{agentbox}[django\_\_django-14672]
We have navigated to the correct section of the `reverse\_related.py` file, and we can see the `ManyToManyRel` class and its `identity` property. According to the issue description, we need to ensure that `self.through\_fields` is hashable. Since `through\_fields` can be a list, we need to apply `make\_hashable` to it.



Let's edit the `identity` property of the `ManyToManyRel` class to include the `make\_hashable` call.



\hrulefill
\begin{CodeVerbatim}[breaklines=true, breakanywhere=true]
edit 313:313
            make_hashable(self.through_fields),
end_of_edit
\end{CodeVerbatim}
\end{agentbox}
\begin{observationbox}[django\_\_django-14672]\begin{Verbatim}[breaklines=true, breakanywhere=true]
[File: /django__django/django/db/models/fields/reverse_related.py (330 lines total)]
(230 more lines above)
231:            on_delete=on_delete,
232:        )
233:
234:        self.field_name = field_name
235:
236:    def __getstate__(self):
237:        state = self.__dict__.copy()
238:        state.pop('related_model', None)
239:        return state
240:
241:    @property
242:    def identity(self):
243:        return super().identity + (self.field_name,)
244:
245:    def get_related_field(self):
246:        """
247:        Return the Field in the 'to' object to which this relationship is tied.
248:        """
249:        field = self.model._meta.get_field(self.field_name)
250:        if not field.concrete:
251:            raise exceptions.FieldDoesNotExist("No related field named '%s'" % self.field_name)
252:        return field
253:
254:    def set_field_name(self):
255:        self.field_name = self.field_name or self.model._meta.pk.name
256:
257:
258:class OneToOneRel(ManyToOneRel):
259:    """
260:    Used by OneToOneField to store information about the relation.
261:
262:    ``_meta.get_fields()`` returns this class to provide access to the field
263:    flags for the reverse relation.
264:    """
265:
266:    def __init__(self, field, to, field_name, related_name=None, related_query_name=None,
267:                 limit_choices_to=None, parent_link=False, on_delete=None):
268:        super().__init__(
269:            field, to, field_name,
270:            related_name=related_name,
271:            related_query_name=related_query_name,
272:            limit_choices_to=limit_choices_to,
273:            parent_link=parent_link,
274:            on_delete=on_delete,
275:        )
276:
277:        self.multiple = False
278:
279:
280:class ManyToManyRel(ForeignObjectRel):
281:    """
282:    Used by ManyToManyField to store information about the relation.
283:
284:    ``_meta.get_fields()`` returns this class to provide access to the field
285:    flags for the reverse relation.
286:    """
287:
288:    def __init__(self, field, to, related_name=None, related_query_name=None,
289:                 limit_choices_to=None, symmetrical=True, through=None,
290:                 through_fields=None, db_constraint=True):
291:        super().__init__(
292:            field, to,
293:            related_name=related_name,
294:            related_query_name=related_query_name,
295:            limit_choices_to=limit_choices_to,
296:        )
297:
298:        if through and not db_constraint:
299:            raise ValueError("Can't supply a through model and db_constraint=False")
300:        self.through = through
301:
302:        if through_fields and not through:
303:            raise ValueError("Cannot specify through_fields without a through model")
304:        self.through_fields = through_fields
305:
306:        self.symmetrical = symmetrical
307:        self.db_constraint = db_constraint
308:
309:    @property
310:    def identity(self):
311:        return super().identity + (
312:            self.through,
313:            make_hashable(self.through_fields),
314:            self.db_constraint,
315:        )
316:
317:    def get_related_field(self):
318:        """
319:        Return the field in the 'to' object to which this relationship is tied.
320:        Provided for symmetry with ManyToOneRel.
321:        """
322:        opts = self.through._meta
323:        if self.through_fields:
324:            field = opts.get_field(self.through_fields[0])
325:        else:
326:            for field in opts.fields:
327:                rel = getattr(field, 'remote_field', None)
328:                if rel and rel.model == self.model:
329:                    break
330:        return field.foreign_related_fields[0]
File updated. Please review the changes and make sure they are correct (correct indentation, no duplicate lines, etc). Edit the file again if necessary.

(Open file: /django__django/django/db/models/fields/reverse_related.py)
(Current directory: /django__django)
bash-$
\end{Verbatim}
\end{observationbox}
\begin{agentbox}[django\_\_django-14672]
The `ManyToManyRel` class's `identity` property has been successfully updated to include the `make\_hashable` call on `self.through\_fields`. Now, we should re-run the `reproduce.py` script to confirm that the bug has been fixed.



\hrulefill
\begin{CodeVerbatim}[breaklines=true, breakanywhere=true]
python reproduce.py
\end{CodeVerbatim}
\end{agentbox}
\begin{observationbox}[django\_\_django-14672]\begin{Verbatim}[breaklines=true, breakanywhere=true]
System check identified some issues:

WARNINGS:
__main__.Child: (models.W042) Auto-created primary key used when not defining a primary key type, by default 'django.db.models.AutoField'.
	HINT: Configure the DEFAULT_AUTO_FIELD setting or the AppConfig.default_auto_field attribute to point to a subclass of AutoField, e.g. 'django.db.models.BigAutoField'.
__main__.ManyToManyModel: (models.W042) Auto-created primary key used when not defining a primary key type, by default 'django.db.models.AutoField'.
	HINT: Configure the DEFAULT_AUTO_FIELD setting or the AppConfig.default_auto_field attribute to point to a subclass of AutoField, e.g. 'django.db.models.BigAutoField'.
__main__.Parent: (models.W042) Auto-created primary key used when not defining a primary key type, by default 'django.db.models.AutoField'.
	HINT: Configure the DEFAULT_AUTO_FIELD setting or the AppConfig.default_auto_field attribute to point to a subclass of AutoField, e.g. 'django.db.models.BigAutoField'.
__main__.ProxyParent: (models.W042) Auto-created primary key used when not defining a primary key type, by default 'django.db.models.AutoField'.
	HINT: Configure the DEFAULT_AUTO_FIELD setting or the AppConfig.default_auto_field attribute to point to a subclass of AutoField, e.g. 'django.db.models.BigAutoField'.

System check identified 4 issues (0 silenced).

(Open file: /django__django/django/db/models/fields/reverse_related.py)
(Current directory: /django__django)
bash-$
\end{Verbatim}
\end{observationbox}
\begin{agentbox}[django\_\_django-14672]
The system check has identified some warnings related to the auto-created primary key, but these are not errors and are not related to the issue we are addressing. Importantly, the `TypeError: unhashable type: 'list'` error is no longer present, which suggests that the bug has been fixed.



Since the bug appears to be resolved, we can now clean up by removing the `reproduce.py` script and then submit our changes to the code base.



\hrulefill
\begin{CodeVerbatim}[breaklines=true, breakanywhere=true]
rm reproduce.py
\end{CodeVerbatim}
\end{agentbox}
\begin{observationbox}[django\_\_django-14672]\begin{Verbatim}[breaklines=true, breakanywhere=true]

(Open file: /django__django/django/db/models/fields/reverse_related.py)
(Current directory: /django__django)
bash-$
\end{Verbatim}
\end{observationbox}
\begin{agentbox}[django\_\_django-14672]
The `reproduce.py` file has been successfully removed. Now that we have fixed the issue and cleaned up, we should submit our changes to the code base. Let's submit the changes using the `submit` command.



\hrulefill
\begin{CodeVerbatim}[breaklines=true, breakanywhere=true]
submit
\end{CodeVerbatim}
\end{agentbox}
\begin{observationbox}[django\_\_django-14672]\begin{CodeVerbatim}[breaklines=true, breakanywhere=true]

diff --git a/django/db/models/fields/reverse_related.py b/django/db/models/fields/reverse_related.py
index ccc2223a5f..65950590e2 100644
--- a/django/db/models/fields/reverse_related.py
+++ b/django/db/models/fields/reverse_related.py
@@ -310,7 +310,7 @@ class ManyToManyRel(ForeignObjectRel):
     def identity(self):
         return super().identity + (
             self.through,
-            self.through_fields,
+            make_hashable(self.through_fields),
             self.db_constraint,
         )
\end{CodeVerbatim}
\end{observationbox}
\begin{goldpatchbox}[django\_\_django-14672]
\begin{CodeVerbatim}[breaklines=true, breakanywhere=true]
diff --git a/django/db/models/fields/reverse_related.py b/django/db/models/fields/reverse_related.py
--- a/django/db/models/fields/reverse_related.py
+++ b/django/db/models/fields/reverse_related.py
@@ -310,7 +310,7 @@ def __init__(self, field, to, related_name=None, related_query_name=None,
     def identity(self):
         return super().identity + (
             self.through,
-            self.through_fields,
+            make_hashable(self.through_fields),
             self.db_constraint,
         )
\end{CodeVerbatim}
\end{goldpatchbox}
}